\chapter{Gestión del tiempo}
\label{cap:gestion-tiempo}

Hasta la etapa 4, el desarrollo se organizó siguiendo el avance de la
materia. Primero se asistía a las clases y se realizaban los trabajos
prácticos para asentar los conocimientos. Luego se comenzaba con la
programación de la etapa correspondiente. Este proceso se complementaba
con la lectura de libros de la bibliografía de la asignatura.

Según las necesidades de cada etapa, primero se elaboraba el esquema de
traducción (EDT) o se realizaba el diseño de la etapa, y luego se pasaba
a la implementación. El desarrollo se hacía de forma incremental: se
implementaban y probaban primero los casos más simples y, a medida que
estos funcionaban correctamente, se avanzaba hacia situaciones más
complejas.

Una vez finalizado el cursado, se terminó la etapa 4 y se comenzó con la
última etapa, correspondiente a la generación de código, y con la
elaboración de este manual.
